\documentclass[12pt,a4paper]{article}
\usepackage[bahasa]{babel}
\usepackage[utf8]{inputenc}
\usepackage{geometry}
\usepackage{graphicx}
\usepackage{hyperref}
\usepackage{listings}
\usepackage{xcolor}
\usepackage{titlesec}
\usepackage{booktabs}
\usepackage{array}
\usepackage{longtable}
\usepackage{fancyhdr}
\usepackage{tikz}
\usetikzlibrary{shapes,arrows,positioning,fit}

\geometry{margin=2.5cm}
\hypersetup{colorlinks=true,linkcolor=blue,urlcolor=cyan}
\pagestyle{fancy}
\fancyhf{}
\rhead{Analisis Sistem Digital Wedding Invitation SaaS}
\lhead{\thepage}
\renewcommand{\headrulewidth}{0.4pt}

\lstset{
  basicstyle=\ttfamily\footnotesize,
  backgroundcolor=\color{gray!10},
  frame=single, breaklines=true,
  keywordstyle=\color{blue}\bfseries,
  commentstyle=\color{green!60!black},
  stringstyle=\color{orange},
  numbers=left, numberstyle=\tiny\color{gray},
  tabsize=2
}

\title{
  \textbf{Analisis Sistem Komprehensif} \\
  \large Digital Wedding Invitation SaaS \\
  \normalsize Next.js App Router + Prisma + Supabase + Zustand
}
\author{Analisis Teknis Arsitektur}
\date{April 2026}

\begin{document}
\maketitle
\tableofcontents
\newpage

% ============================================================
\section{Gambaran Umum Sistem}
% ============================================================

Sistem ini adalah \textbf{platform SaaS undangan pernikahan digital} berbasis web yang
memungkinkan pengguna membuat, mengedit, dan mempublikasikan undangan pernikahan
dengan berbagai template premium. Sistem dibangun menggunakan:

\begin{table}[h]
\centering
\begin{tabular}{>{\bfseries}l l l}
\toprule
Komponen & Teknologi & Fungsi \\
\midrule
Framework      & Next.js 15 (App Router) & SSR + RSC + Routing \\
Database ORM   & Prisma 7 + PostgreSQL   & Persistensi data \\
Auth Provider  & Supabase Auth           & Autentikasi JWT \\
State Mgmt     & Zustand                 & Editor client state \\
Storage        & Supabase Storage        & File/Gambar/Musik \\
UI Animation   & Framer Motion           & Template animations \\
Styling        & Tailwind CSS            & Utility-first CSS \\
\bottomrule
\end{tabular}
\caption{Stack Teknologi Utama}
\end{table}

% ============================================================
\section{Arsitektur Aplikasi}
% ============================================================

\subsection{Struktur Direktori}

\begin{lstlisting}[language=bash, caption=Struktur Direktori Utama]
src/
  app/
    (admin)/          # Route group admin
    (auth)/           # Login / Register
    (create)/create/  # Editor undangan
    (invitation)/     # Public view undangan
    (preview)/        # Preview template
    (user)/           # Dashboard user
    api/              # API endpoints
  components/
    templates/        # 11+ template komponen
    user/create/      # Editor UI components
    ui/               # Shared UI (shadcn)
  lib/
    actions/          # Server Actions (Next.js)
    templateRegistry.ts
    prisma.ts
  store/
    useEditorStore.ts # Zustand global state
  types/
    invitation.ts     # TypeScript interfaces
\end{lstlisting}

\subsection{Route Group Architecture}

Next.js App Router menggunakan \textit{Route Groups} dengan tanda kurung
untuk pengelompokan tanpa mempengaruhi URL. Sistem ini memiliki 6 route group:

\begin{enumerate}
  \item \texttt{(admin)} --- Panel administrasi (user, token, template CRUD)
  \item \texttt{(auth)} --- Halaman login dan registrasi via Supabase
  \item \texttt{(create)} --- Editor undangan dengan live preview
  \item \texttt{(invitation)} --- Halaman publik undangan \texttt{/undangan/[slug]}
  \item \texttt{(preview)} --- Preview template \texttt{/preview/[template]}
  \item \texttt{(user)} --- Dashboard pengguna
\end{enumerate}

% ============================================================
\section{Skema Database (Prisma)}
% ============================================================

\subsection{Model Data}

Database menggunakan PostgreSQL dengan Prisma sebagai ORM. Berikut adalah
relasi antar model:

\begin{table}[h]
\centering
\begin{tabular}{>{\bfseries}l l l}
\toprule
Model & Field Kunci & Relasi \\
\midrule
User       & id, email, tokens, role, status   & has many Invitation \\
Template   & id, title, category, configPath    & has many Invitation \\
Invitation & id, slug, contentData (JSON), userId, templateId & belongs to User, Template \\
GuestWish  & id, guestName, message, isPresent  & belongs to Invitation \\
Music      & id, title, url, artist             & standalone \\
\bottomrule
\end{tabular}
\caption{Relasi Antar Model Database}
\end{table}

\subsection{Field contentData (JSON)}

Field \texttt{contentData} bertipe \texttt{Json?} di Prisma dan di-cast ke
interface TypeScript \texttt{InvitationContent}:

\begin{lstlisting}[language=typescript, caption=Interface InvitationContent]
interface CoupleDetail {
  nama: string;         // Nama mempelai
  ortu_ayah: string;    // Nama ayah
  ortu_ibu: string;     // Nama ibu
  instagram?: string;
  foto?: string;        // URL foto
}

interface EventDetail {
  tipe: string;         // "Akad Nikah" | "Resepsi"
  tanggal: string;
  jam: string;
  lokasi: string;
  alamat_lengkap: string;
  link_maps: string;
}

interface InvitationContent {
  mempelai_pria: CoupleDetail;
  mempelai_wanita: CoupleDetail;
  acara: EventDetail[];
  love_story: LoveStory[];
  gallery: string[];
  digital_envelope: DigitalEnvelope[];
  guest_wishes: GuestWishEntry[];
  music_url?: string;
  rsvp_url?: string;
  countdown_date?: string;
  cover_title?: string;
  hero_image?: string;
  wedding_date?: string;
  dress_code?: string;
  closing_message?: string;
}
\end{lstlisting}

% ============================================================
\section{Sistem Autentikasi}
% ============================================================

Autentikasi menggunakan \textbf{Supabase Auth} dengan JWT session.
Alur autentikasi adalah sebagai berikut:

\begin{enumerate}
  \item User melakukan login via Supabase UI
  \item Supabase mengembalikan JWT token yang disimpan sebagai cookie
  \item Setiap Server Component/Action memanggil \texttt{createServerSupabase()} untuk memvalidasi sesi
  \item Jika session valid, \texttt{user.id} digunakan untuk query Prisma
  \item Role \texttt{admin} memiliki akses ke route \texttt{(admin)}
\end{enumerate}

\subsection{Token Economy}

Sistem menggunakan \textit{token economy} unik:
\begin{itemize}
  \item Setiap user memiliki field \texttt{tokens: Int}
  \item Membuat undangan baru membutuhkan \textbf{10 token}
  \item Token dikurangi secara atomik via Prisma transaction
  \item Admin dapat menambah token via dashboard admin
\end{itemize}

% ============================================================
\section{Sistem Template}
% ============================================================

\subsection{Arsitektur Template Registry}

Template di-register secara statis di \texttt{src/lib/templateRegistry.ts}.
Sistem key normalisasi:

\begin{lstlisting}[language=typescript, caption=Key Normalization]
// Title template -> key registry
// "Gen Z Pastel" -> "genzpastel"
const key = title.replace(/\s+/g, "").toLowerCase();

export const templateRegistry = {
  nerogold: NeroGold,
  auradark: AuraDark,
  genzpastel: GenZPastel,
  // ...11 template total
};
\end{lstlisting}

\subsection{Template yang Tersedia}

\begin{table}[h]
\centering
\begin{tabular}{>{\bfseries}l l l}
\toprule
Key & Nama Template & Kategori \\
\midrule
nerogold          & Nero Gold            & Modern \\
auradark          & Aura Dark            & Modern \\
pink              & Pink                 & Romantis \\
royal             & Royal                & Klasik \\
jawaroyalkeraton  & Jawa Royal Keraton   & Tradisional \\
sundaanggunpriangan & Sunda Anggun Priangan & Tradisional \\
minangmaharaja    & Minang Maharaja      & Tradisional \\
balisacredluxury  & Bali Sacred Luxury   & Tradisional \\
batakheritage     & Batak Heritage       & Tradisional \\
bugisgoldensilk   & Bugis Golden Silk    & Tradisional \\
genzpastel        & Gen Z Pastel         & Gen-Z \\
\bottomrule
\end{tabular}
\caption{Daftar Template Terdaftar}
\end{table}

\subsection{Anatomi Template}

Setiap template adalah React Client Component dengan struktur:

\begin{lstlisting}[language=typescript, caption=Struktur Template Wajib]
"use client"; // WAJIB baris pertama

import { InvitationContent } from "@/types/invitation";

export default function NamaTemplate({
  data
}: {
  data: InvitationContent
}) {
  // useState untuk: opened, playing, countdown
  // useEffect untuk: audio, countdown timer, auto-slider
  return ( /* JSX */ );
}
\end{lstlisting}

\subsection{Alur Resolusi Template}

\begin{enumerate}
  \item Database menyimpan \texttt{Invitation.templateId} (relasi ke \texttt{Template.id})
  \item Saat halaman publik \texttt{/undangan/[slug]} diakses, Server Component fetch undangan beserta template
  \item \texttt{template.title} dinormalisasi: \texttt{title.replace(/\textbackslash s+/g,"").toLowerCase()}
  \item Key digunakan untuk lookup di \texttt{templateRegistry}
  \item Komponen yang ditemukan di-render dengan \texttt{data={contentData}}
\end{enumerate}

% ============================================================
\section{Sistem Editor (Live Preview)}
% ============================================================

\subsection{Komponen Editor}

Editor terdiri dari 4 komponen utama:

\begin{table}[h]
\centering
\begin{tabular}{>{\bfseries}l l}
\toprule
Komponen & Fungsi \\
\midrule
\texttt{CreateDesktopView} & Layout utama grid 2 kolom \\
\texttt{EditorSidebar}     & Form input semua field + template picker \\
\texttt{LivePreview}       & Preview real-time via MobileDeviceEmulator \\
\texttt{SectionCard}       & Accordion UI per section \\
\bottomrule
\end{tabular}
\caption{Komponen Editor}
\end{table}

\subsection{State Management dengan Zustand}

\begin{lstlisting}[language=typescript, caption=useEditorStore State]
interface EditorState {
  formData: InvitationContent;  // Data undangan
  invitationId: string;         // ID undangan aktif
  activeSection: SectionId;     // Section form aktif
  activeTemplate: string;       // Title template aktif
  isSaving: boolean;            // Status auto-save

  setFormData: (data: Partial<InvitationContent>) => void;
  // Merge partial update dengan spread operator
}
\end{lstlisting}

\subsection{Alur Data Editor}

\begin{enumerate}
  \item \texttt{CreateDesktopView} (Server Component) fetch data awal dari Prisma
  \item Data dikirim ke \texttt{EditInvitationClient} sebagai props
  \item \texttt{useEffect} menginisialisasi Zustand store dengan data awal
  \item User mengetik di \texttt{EditorSidebar} → \texttt{setFormData()} dipanggil → store update
  \item \texttt{LivePreview} subscribe ke \texttt{useEditorStore} → re-render template
  \item \texttt{useDebouncedCallback} (1500ms) trigger \texttt{updateInvitationContent()} Server Action
  \item Server Action validasi ownership → update Prisma → \texttt{revalidatePath()}
\end{enumerate}

\subsection{Fix: Mencegah Template Remount}

\textbf{Masalah kritis yang ditemukan:} Sebelumnya \texttt{LivePreview} memanggil
\texttt{renderTemplate()} sebagai fungsi biasa, menyebabkan React menganggap
\texttt{Component} sebagai instance baru setiap render → template remount →
semua \texttt{useState} (termasuk \texttt{opened}) reset.

\textbf{Solusi:}

\begin{lstlisting}[language=typescript, caption=Fix useMemo di LivePreview]
// SEBELUM (buggy): Component berubah reference tiap render
const renderTemplate = () => {
  const Component = getTemplate[key]; // baru tiap render!
  return <Component data={formData} />;
};

// SESUDAH (fixed): Component hanya berubah jika template berganti
const Component = useMemo(() => {
  return getTemplate[normalizeTemplateName(activeTemplate)];
}, [activeTemplate]); // formData tidak masuk dependency
\end{lstlisting}

% ============================================================
\section{Alur Publik Undangan}
% ============================================================

\subsection{Public Page Flow}

\begin{enumerate}
  \item Tamu mengakses \texttt{/undangan/[slug]}
  \item Next.js Server Component fetch: \texttt{prisma.invitation.findUnique(\{where:\{slug\}, include:\{template:true\}\})}
  \item \texttt{contentData} di-cast ke \texttt{InvitationContent}
  \item Template dinormalisasi dan di-lookup di registry
  \item \texttt{PublicInvitationClient} di-render dengan data
  \item Template client-side menghandle: cover animation, musik autoplay, countdown, RSVP
\end{enumerate}

\subsection{SEO dan Caching}

\begin{itemize}
  \item Halaman public di-render sebagai \textbf{RSC (React Server Component)} untuk SEO
  \item \texttt{revalidatePath()} dipanggil setelah setiap update untuk invalidasi ISR cache
  \item Metadata halaman dapat di-generate dari \texttt{contentData}
\end{itemize}

% ============================================================
\section{Server Actions}
% ============================================================

Sistem menggunakan \textbf{Next.js Server Actions} untuk operasi database,
menggantikan API Routes tradisional:

\begin{table}[h]
\centering
\begin{tabular}{>{\ttfamily\footnotesize}l l l}
\toprule
\textrm{Action} & Fungsi & Auth \\
\midrule
createInitialInvitation & Buat undangan baru (deduct 10 token) & Required \\
updateInvitationContent & Update content + template            & Required + Owner \\
deleteInvitation        & Hapus undangan                       & Required + Owner \\
getInvitationById       & Fetch undangan by ID                 & Public \\
saveInvitation          & Quick save tanpa auth check          & None \\
\bottomrule
\end{tabular}
\caption{Daftar Server Actions}
\end{table}

\subsection{Keamanan Server Actions}

\begin{itemize}
  \item Setiap action destructive memvalidasi \texttt{user.id} dari Supabase session
  \item Ownership check: \texttt{invitation.userId === user.id}
  \item Prisma transaction digunakan untuk operasi atomik (token deduction + invitation create)
  \item Error P2002 (unique constraint) ditangani secara eksplisit untuk slug conflicts
\end{itemize}

% ============================================================
\section{Arsitektur Gen Z Pastel Template}
% ============================================================

\subsection{Design System}

Template \textit{Gen Z Pastel} mengimplementasikan design system berikut:

\begin{table}[h]
\centering
\begin{tabular}{>{\bfseries}l l l}
\toprule
Token & Nilai & Penggunaan \\
\midrule
bg      & \texttt{\#FDFBF7} & Background utama \\
beige   & \texttt{\#E5E0D8} & Card background \\
sage    & \texttt{\#A3B19B} & Accent + button + border \\
sageDark & \texttt{\#7A9070} & Heading accent \\
text    & \texttt{\#4A4A4A} & Body text \\
muted   & \texttt{\#888880} & Secondary text \\
\bottomrule
\end{tabular}
\caption{Color Tokens Gen Z Pastel}
\end{table}

\subsection{Seksi Template (Berurutan)}

\begin{enumerate}
  \item \textbf{Cover} --- Full-screen overlay, decorative blobs, pulse button animation
  \item \textbf{Hero} --- Parallax background, nama mempelai besar dengan italic serif
  \item \textbf{Mempelai} --- Arch-shaped photo (CSS border-radius asimetris), info orang tua
  \item \textbf{Countdown} --- Real-time timer dengan \texttt{setInterval}, 4 kolom grid
  \item \textbf{Jadwal Acara} --- Vertical timeline, map CTA link
  \item \textbf{Love Story} --- Auto-slider 4.8 detik, progress pill dots
  \item \textbf{Gallery} --- CSS columns-2 masonry, aspect-ratio bervariasi
  \item \textbf{Amplop Digital} --- Copy to clipboard, QRIS image support
  \item \textbf{RSVP} --- CTA button link + sticky footer CTA
  \item \textbf{Ucapan} --- Infinite horizontal marquee animation
  \item \textbf{Closing} --- Quote penutup, nama mempelai serif
\end{enumerate}

\subsection{State Lokal Template}

\begin{lstlisting}[language=typescript, caption=State Template]
const [opened, setOpened]    = useState(false);  // Cover toggle
const [playing, setPlaying]  = useState(true);   // Audio state
const [story, setStory]      = useState(0);       // Love story index
const [copied, setCopied]    = useState({});      // Copy feedback
const [tl, setTl]            = useState({d,h,m,s}); // Countdown
const audio                  = useRef<Audio>();   // Audio ref
\end{lstlisting}

% ============================================================
\section{Seed System}
% ============================================================

Template didaftarkan ke database via script \texttt{seed-templates.ts}:

\begin{lstlisting}[language=bash, caption=Menjalankan Seed]
pnpm seed:templates
# => dotenv --file .env.local -- tsx seed-templates.ts
\end{lstlisting}

Logika seed menggunakan \textit{upsert pattern} — jika template dengan
title yang sama sudah ada, dilewati (\texttt{[DILEWATI]}), jika belum
ada maka dibuat (\texttt{[DIBUAT]}).

% ============================================================
\section{Diagram Alur Sistem}
% ============================================================

\subsection{Alur Pembuatan Undangan}

\begin{center}
\begin{tikzpicture}[
  node distance=1.2cm,
  box/.style={rectangle, rounded corners, draw, minimum width=3.5cm, minimum height=0.8cm, font=\small},
  arrow/.style={->, thick}
]
\node[box, fill=blue!10]   (login)    {User Login (Supabase)};
\node[box, fill=green!10]  (dash)     [below=of login]    {Dashboard /dashboard};
\node[box, fill=yellow!10] (create)   [below=of dash]     {Klik "Buat Undangan"};
\node[box, fill=orange!10] (token)    [below=of create]   {Cek Token >= 10};
\node[box, fill=red!10]    (db)       [below=of token]    {Prisma Transaction};
\node[box, fill=purple!10] (editor)   [below=of db]       {Editor /create};
\node[box, fill=teal!10]   (preview)  [below=of editor]   {Live Preview};
\node[box, fill=gray!10]   (save)     [below=of preview]  {Auto-save (1500ms debounce)};
\node[box, fill=blue!20]   (public)   [below=of save]     {Publik /undangan/[slug]};

\draw[arrow] (login) -- (dash);
\draw[arrow] (dash)  -- (create);
\draw[arrow] (create)-- (token);
\draw[arrow] (token) -- node[right]{\tiny -10 token} (db);
\draw[arrow] (db)    -- (editor);
\draw[arrow] (editor)-- (preview);
\draw[arrow] (preview)-- (save);
\draw[arrow] (save)  -- (public);
\end{tikzpicture}
\end{center}

\subsection{Alur Resolusi Template di Halaman Publik}

\begin{center}
\begin{tikzpicture}[
  node distance=1.1cm,
  box/.style={rectangle, rounded corners, draw, minimum width=4.5cm, minimum height=0.75cm, font=\small},
  arrow/.style={->, thick}
]
\node[box, fill=blue!10]   (url)    {GET /undangan/[slug]};
\node[box, fill=green!10]  (fetch)  [below=of url]  {Prisma: findUnique(slug) + include template};
\node[box, fill=yellow!10] (norm)   [below=of fetch] {Normalize: title.replace(\textbackslash s+,"").toLowerCase()};
\node[box, fill=orange!10] (reg)    [below=of norm]  {Lookup templateRegistry[key]};
\node[box, fill=purple!10] (comp)   [below=of reg]   {<Component data=\{contentData\} />};
\node[box, fill=teal!10]   (render) [below=of comp]  {Client Render: animasi, musik, countdown};
\draw[arrow] (url)   -- (fetch);
\draw[arrow] (fetch) -- (norm);
\draw[arrow] (norm)  -- (reg);
\draw[arrow] (reg)   -- (comp);
\draw[arrow] (comp)  -- (render);
\end{tikzpicture}
\end{center}

% ============================================================
\section{Temuan dan Rekomendasi}
% ============================================================

\subsection{Temuan Teknis}

\begin{enumerate}
  \item \textbf{Bug LivePreview (Fixed)} --- \texttt{Component} reference baru setiap render menyebabkan remount template. Diselesaikan dengan \texttt{useMemo([activeTemplate])}.
  \item \textbf{Schema Mismatch} --- Template awal menggunakan field yang tidak sesuai dengan \texttt{InvitationContent} (e.g., \texttt{nama\_lengkap} vs \texttt{nama}). Diselesaikan dengan menyesuaikan ke schema resmi.
  \item \textbf{Audio Management} --- \texttt{useRef} digunakan untuk audio agar tidak trigger re-render saat play/pause.
  \item \textbf{Token Race Condition} --- Ditangani dengan Prisma \texttt{\$transaction} dan atomik \texttt{decrement}.
\end{enumerate}

\subsection{Rekomendasi Pengembangan}

\begin{itemize}
  \item Tambahkan \textbf{image optimization} via Supabase Storage transform API untuk gallery
  \item Implementasi \textbf{WebSocket} atau \textbf{Server-Sent Events} untuk collaborative editing
  \item Tambahkan \textbf{Zod validation} di Server Actions untuk type-safe input
  \item Pertimbangkan \textbf{Redis caching} untuk halaman undangan yang sering dikunjungi
  \item Tambahkan \textbf{OG Image generation} (via \texttt{@vercel/og}) per undangan untuk WhatsApp preview
  \item Implementasi \textbf{A/B testing} template dengan analytics konversi
\end{itemize}

% ============================================================
\section{Kesimpulan}
% ============================================================

Sistem Digital Wedding Invitation SaaS ini merupakan implementasi yang
solid dari \textit{full-stack modern web architecture} dengan beberapa keunggulan:

\begin{itemize}
  \item \textbf{Separation of Concerns} yang baik antara template, editor, dan data layer
  \item \textbf{Type Safety} end-to-end dengan TypeScript + Prisma + Zod
  \item \textbf{Mobile-First} template design yang konsisten
  \item \textbf{Scalable Template System} via registry pattern
  \item \textbf{Secure Server Actions} dengan ownership validation
  \item \textbf{Real-time Preview} dengan debounced auto-save yang efisien
\end{itemize}

Sistem siap untuk production dengan beberapa improvement yang direkomendasikan
di atas untuk meningkatkan performa dan pengalaman pengguna.

\vspace{1cm}
\hrule
\vspace{0.3cm}
\noindent\small\textit{Dokumen ini dihasilkan dari analisis kode sumber project
\texttt{d:\textbackslash memekgosong\textbackslash wedding} pada April 2026.}

\end{document}
